\section{Fake it, triangulação e implementação óbvia}

Para chegar ao verde há dois caminhos, ambos de Kent Beck
(\citetitle{livro:beck2002}):

\begin{description}
\item[Sem saber a solução.] Fingir (\emph{fake it}): devolver a constante
  que o teste espera, por exemplo \texttt{return 100.0}. Em seguida
  \emph{triangular}: escrever um segundo exemplo que a constante não
  satisfaz, o que força a generalizar.
\item[Sabendo a solução.] Escrever direto a \emph{implementação óbvia}.
\end{description}

Ficou vermelho sem entender por quê? O passo foi grande demais. Volte,
desfaça, e dê um passo menor.
